\section{Introduction}
\subsection{Super Thrall's problem}
Thrall \cite{Thrall} famously introduced a certain canonical decomposition of the tensor algebra coming from free Lie algebras:
\begin{equation}\label{eq:Thrall}
    \Trm(V) \cong \bigoplus_\lambda \Lcal_\lambda(V).
\end{equation}
Here $V = \CC^N$, $\lambda$ ranges over all integer partitions, the $\Lcal_\lambda = \Lcal_\lambda(V)$ are called \textit{Lie modules}, and the isomorphism is as $\GL(V)$-modules. See \Cref{sec:Thrall-classical} for details. \textit{Thrall's problem} is to explicitly decompose the $\Lcal_\lambda$ in terms of irreducible representations. Equivalently, we seek the Schur decomposition of the characters $\ch(\Lcal_\lambda; \xb)$ in the limit $N \to \infty$, where $\xb = x_1, x_2, \ldots$. Thrall's problem remains open in general, and has received significant attention in the literature since its formulation \cite{AS,BBG,Brandt,Gar,GR,Kly,KW,Sch,Sun}. See Vic Reiner's lecture slides \cite{Rei} for a lucid introduction to Thrall's problem and Adriano Garsia's lecture notes \cite{Gar} for a detailed introduction to the theory of free Lie algebras and many of the results we extend.

We study a supersymmetric extension of Thrall's problem. In \Cref{thm:super-Thrall}, we extend \eqref{eq:Thrall} to a canonical decomposition of the tensor superalgebra coming from free Lie superalgebras:
\begin{equation}\label{eq:super-Thrall}
    \Trmtil(V) \cong \bigoplus_A \widetilde{\Lcal}_A(V).
\end{equation}
Here $V = V_0 \oplus V_1 = \CC^N \oplus \CC^M$ is a $\ZZ/2$-graded vector space (a \textit{super vector space}), $A$ ranges over all infinite $\ZZ_{\geq 0}$-valued matrices $A=(a_{i,j})_{i,j \geq 0}$ with finite support and $a_{0,0}=0$, the $\Ltil_A = \Ltil_A(V)$ are called \textit{super Lie modules}, and the isomorphism is as $\GL(\CC^N) \oplus \GL(\CC^M)$-modules. See \Cref{sec:super-Thrall} for details. We consider the problem of explicitly decomposing $\Ltil_A$ into irreducible representations when $N=M \to \infty$ under the diagonal action $\GL(\CC^N) \hookrightarrow \GL(\CC^N) \oplus \GL(\CC^N)$. Equivalently, we seek the Schur decomposition of the characters $\ch(\Ltil_A; \xb, \xb)$. We recover the classical case as $\Ltil_A = \Lcal_\lambda$ when $a_{i,0}$ is the number of copies of $i$ in $\lambda$, and $a_{i,j} = 0$ for $j > 0$.

\subsection{A key special case}
The most important component of \eqref{eq:Thrall} is $\lambda = (n)$, where $\Lcal_n = \Lcal_{(n)}$ is the degree-$n$ homogeneous component of the free Lie algebra of $V$ in $\Trm(V)$. The decomposition of the higher Lie modules $\Lcal_\lambda$ may be determined by $\ch(\Lcal_n)$ from the Littlewood--Richardson rule and plethysm.

In the super case, $\widetilde{\Lcal}_{n,m}$ is the bidegree-$(n,m)$ homogeneous component of the free Lie superalgebra of $V_0 \oplus V_1$ in $\Trmtil(V_0 \oplus V_1)$. We have $\widetilde{\Lcal}_{n,0} = \Lcal_n$ and $\widetilde{\Lcal}_{n,m} = \widetilde{\Lcal}_A$ where $a_{n,m}=1$ and $a_{i,j} = 0$ otherwise. As before, the decomposition of the higher super Lie modules $\widetilde{\Lcal}_A$ may be determined by $\ch(\widetilde{\Lcal}_{n,m})$ from the Littlewood--Richardson rule and plethysm.

\textit{Brandt's formula} \cite{Brandt} gives the decomposition of $\Lcal_n$ in the power sum basis. We prove the following generalization:
\begin{thm}\label{thm:super-brandt}
    The $\GL(\CC^N)$-character of $\Ltil_{n,m}$ is given by
    \begin{equation}
        \ch(\Ltil_{n,m}; \xb_N) = \frac{1}{n+m} \sum_{d \mid \gcd(n,m)} (-1)^{m + \frac{m}{d}} \mu(d) \binom{\frac{n+m}{d}}{\frac{m}{d}} p_d(\xb_N)^\frac{n+m}{d}.
    \end{equation}
\end{thm}
The $m=0$ case of \Cref{thm:super-brandt} reduces to Brandt's formula. We also obtain the dimension of $\Ltil_{n, m}$, generalizing \textit{Witt's formula} \cite{Witt} in the $m=0$ case:

\begin{coro}[{\cite{Pet}}]\label{cor:super-Witt}
    When $V_0 = V_1 = \mathbb{C}^N$ and $n+m > 0$,
      \[ \dim \Ltil_{n,m} = \frac{1}{n+m} \sum_{d \mid \gcd(n, m)} (-1)^{m + \frac{m}{d}} \mu(d) \binom{\frac{n+m}{d}}{\frac{m}{d}} N^{\frac{n+m}{d}}. \]
\end{coro}

Klyachko \cite{Kly} described the Schur--Weyl dual of $\Lcal_n$ in terms of certain induced representations. To state our generalization, we need some notation. Let $C_{n+m}$ be the cyclic subgroup of the symmetric group $S_{n+m}$ generated by the long cycle $\pi_{n+m} = (1\,2\,\cdots\,n+m)$. Let $\chi^r \colon C_{n+m} \to \CC$ be the character with $\chi^r(\pi_{n+m}) = \exp(2\pi ir/(n+m))$. Finally, let $\chi^\mathrm{cyc}$ be the character of the action where $\pi_{n+m}$ cyclically increments $m$-element subsets of $[n+m]$. We prove the following: 
\begin{thm}\label{thm:super-induced}
We have
\begin{equation}
    \Ch(\Ltil_{n,m}) =
    \begin{cases}
        \FrobCh((\chi^\mathrm{cyc} \otimes \chi^1)\uparrow_{C_{n+m}}^{S_{n+m}}) & \text{if $m$ is odd} \\
        \FrobCh((\chi^\mathrm{cyc} \otimes \chi^{m/2+1})\uparrow_{C_{n+m}}^{S_{n+m}}) & \text{if $m$ is even.}
    \end{cases}
\end{equation}
\end{thm}
When $m=0$, the $\chi^\mathrm{cyc}$ factor may be omitted, recovering the $\Lcal_n$ case.

\subsection{Super Schur decompositions}

Kra{\'s}kiewicz--Weyman determined the Schur decomposition of $\Lcal_n$ \cite{KW}. We generalize this result to $\widetilde{\Lcal}_{n,m}$. Our argument introduces several super extensions of existing concepts in tableau combinatorics.

Let $\SYT_\pm(\lambda)$ denote the set of standard Young tableaux of shape $\lambda \vdash n$ where each cell is marked as \textit{positive} or \textit{negative}. For $\mathcal{T} \in \SYT_\pm(\lambda)$, let $\Neg(\mathcal{T}) \subseteq [n]$ be the set of all $i$ such that the box containing $i$ in $\mathcal{T}$ is marked as negative. Write $\negg(\mathcal{T}) \coloneqq |\Neg(\mathcal{T})|$.
\begin{defi}
    The \textit{(super) descent set} of $\mathcal{T} \in \SYT_\pm(\lambda)$ is the set $\Des(\mathcal{T})$ of all $i=1, \ldots, n-1$ such that either $i+1 \not\in \Neg(\mathcal{T})$ and $i+1$ appears in a strictly lower row of $\mathcal{T}$ than $i$, or $i \in \Neg(\mathcal{T})$ and $i+1$ does not appear in a strictly lower row of $\mathcal{T}$ than $i$. The \textit{(super) major index} of $\mathcal{T} \in \SYT_\pm(\lambda)$ is
\begin{equation}\label{eq:super-maj}
    \maj(\mathcal{T}) = \sum_{i \in \Des(\mathcal{T})} i.
\end{equation}
\end{defi}
Here we use English notation, with the longest row on top, and $\Neg(\mathcal{T}) = \varnothing$ recovers the usual descent set and major index. 

Our super generalization of Kra{\'s}kiewicz--Weyman's theorem is as follows. When $m=0$, we recover Kra{\'s}kiewicz--Weyman's original result \cite{KW} in type $A$.
\begin{thm}\label{thm:super-KW}
    The multiplicity of the Schur module $V^\lambda$ in $\Ltil_{n,m}$ is given by
    \begin{equation}\label{eq:super-KW}
        |\{ \mathcal{T} \in \SYT_\pm(\lambda) : \maj(\mathcal{T}) \equiv_{n+m} 1, \negg(\mathcal{T}) = m \}|.
    \end{equation}
\end{thm}

Our proof of \Cref{thm:super-KW} involves super analogues of the theory of $P$-partitions and symmetric function identities. For instance, we prove the following $q,t$-hook formula, which was announced in \cite[(6.1)]{BS}:
\begin{thm}\label{thm:maj-neg-hook}
    For any $\lambda \vdash n$,
    \begin{equation}\label{eq:maj-neg-hook}
        \sum_{\mathcal{T} \in \SYT_\pm(\lambda)} q^{\maj(\mathcal{T})} t^{\negg(\mathcal{T})} = [n]_q! \prod_{(r, c) \in \lambda} \frac{q^{r-1} + tq^{c-1}}{[h(r, c)]_q}
    \end{equation}
    where $[n]_q \coloneqq 1 + q + \cdots + q^{n-1}$ and $[n]_q! \coloneqq [1]_q[2]_q \cdots [n]_q$.
\end{thm}

By contrast, Kra{\'s}kiewicz--Weyman's proof uses a character computation to identify $\Lcal_n$ as isomorphic to a submodule of the classical coinvariant algebra. See \cite[\S 8]{Gar} for a detailed discussion. The coinvariant algebra is very well-studied and intimately related to geometry and topology. It would be interesting to extend this connection, which we leave as an open problem.
\begin{prob}
    Find an interpretation of \Cref{thm:super-KW} which directly relates super Lie modules and coinvariant algebras.
\end{prob}

\subsection{Paper organization}
 The rest of the paper is organized as follows. In \Cref{sec:background} we recall some necessary background on (super)symmetric functions and representation theory, and introduce the classical case of Thrall's problem. In \Cref{sec:Thrall} we define (free) Lie superalgebras and prove the super analogue of Brandt's formula (\Cref{thm:super-brandt}), before explicitly defining Thrall's problem for free Lie superalgebras. In \Cref{sec:super-induced} we determine the Schur--Weyl dual of $\Ltil_{n,m}$, proving \Cref{thm:super-induced}. In \Cref{sec:super-tab} we develop some super tableau combinatorics and prove the $q,t$-hook formula (\Cref{thm:maj-neg-hook}), which is in turn used to prove \Cref{thm:super-KW} in \Cref{sec:super-KW}. Further directions are discussed in \Cref{sec:addl}.

\section{Background}\label{sec:background}
\subsection{Symmetric functions}\label{sec:symf}
We first review some background on tableau combinatorics and symmetric functions, and fix notation. For full details, see \cite[Chapter~7]{Stan}. An \emph{integer partition} $\lambda$ is a weakly decreasing sequence of positive integers $\lambda = (\lambda_1 \ge \cdots \ge \lambda_k > 0)$, and we let $\mathrm{Par}$ denote the set of all integer partitions. If $\sum_i \lambda_i = n$, then we denote this by $\lambda \vdash n$. By a slight abuse of notation we identify $\lambda$ with its \emph{Ferrers diagram}, which is the subset of $\ZZ_{>0} \times \ZZ_{>0}$ consisting of $\lambda_i$ left-justified cells in row $i$. We draw our partitions using English notation, so that the longest row appears at the top. The \emph{length} of $\lambda$, denoted $\ell(\lambda)$, is the number of nonempty rows in its Ferrers diagram. We sometimes will also use the notation $\lambda = (1^{a_1}2^{a_2} \cdots)$ to denote the partition with $a_i$ parts equal to $i$. 

Assume $\lambda \vdash n$, and let $[n] = \{ 1, 2, \ldots, n \}$. A \emph{standard tableau} of shape $\lambda$ is a bijective map $T : \lambda \to [n]$ that is strictly increasing along the rows and columns of $\lambda$. The set of all standard tableaux of shape $\lambda$ will be denoted by $\SYT(\lambda)$. For $i = 1, \ldots, n-1$, we say that $i$ is a \emph{descent} of $T$ if $i+1$ appears in a strictly lower row of $T$ than $i$. Let
\[
\Des(T) \coloneqq \{ i : i \text{ is a descent of } T \} \subseteq [n-1].
\]
